\pdfoutput=1
\documentclass[11pt, a4paper]{article}

% -------------------------------------------------------------------------
% PACKAGES
% -------------------------------------------------------------------------
\usepackage[utf8]{inputenc}
\usepackage{amsmath}
\usepackage{amsfonts, amssymb}
\usepackage{geometry}
\usepackage{microtype}
\usepackage{authblk}
\usepackage{fancyhdr}
\usepackage{parskip} % Ensure that a blank line is automatically added between every paragraph without having to manually type \vspace{1em} each time.
\usepackage{hyperref}
\geometry{margin=1in, headheight=14.5pt}

% -------------------------------------------------------------------------
% HEADER & FOOTER
% -------------------------------------------------------------------------
\pagestyle{fancy}
\fancyhf{}
\lhead{Mechanistic Physics}
\rhead{Marab\={u}t's Theory of Gravity: Manuscript 02 of 19}
\cfoot{\thepage}

% -------------------------------------------------------------------------
% TITLE & AUTHOR
% -------------------------------------------------------------------------
\title{\textbf{Marab\={u}t's Theory of Gravity: \\ Resolving the Weyl Potential Anomaly} \\ \large{Fluid Refraction and Dynamic Superposition in the Electron Flow Model}}

\author[1]{Christopher Berry Marab\={u}t\thanks{ORCID iD: \href{https://orcid.org/0009-0001-9759-9587}{0009-0001-9759-9587} | Email: chris.marabut@nestlink.com}}
\author[1]{Angelina Lopez Marab\={u}t\thanks{ORCID iD: \href{https://orcid.org/0009-0007-9937-6145}{0009-0007-9937-6145} | Email: angelina.marabut@nestlink.com}}

\affil[1]{Independent Researchers}

\date{May 6, 2026 \\[1.5em] \small{DOI: \url{https://doi.org/10.5281/zenodo.20060223} $|$ Version 3}}

\begin{document}

\hypersetup{
  pdftitle={Marab\={u}t's Theory of Gravity: Resolving the Weyl Potential Anomaly},
  pdfauthor={Christopher Berry Marab\={u}t, Angelina Lopez Marab\={u}t},
  pdfkeywords={Weyl potential, dark energy survey, fluid refraction, general relativity anomaly, cosmology, electron flow model, efm, vacuum flux},
  pdfsubject={Astrophysics Brief Communication on the Weyl Potential Anomaly}
}

\maketitle

% -------------------------------------------------------------------------
% ABSTRACT
% -------------------------------------------------------------------------
\begin{abstract}
Recent analyses of the Dark Energy Survey (DES) have revealed a $3\sigma$ discrepancy in the evolution of the Weyl potential, indicating that gravitational wells 3.5 to 5 billion years ago were shallower than predicted by General Relativity. Traditional kinematic models, which rely on the static geometric curvature of spacetime, struggle to account for this without invoking shifting dark energy paradigms. We propose that these findings are a direct macroscopic signature of the Electron Flow Model (EFM). By redefining gravity as hydrodynamic fluid pressure generated by vacuum flux consumption, the EFM replaces the ``deformation of spacetime'' with angular fluid refraction and viscous propagation. Furthermore, the EFM demonstrates that shifting baseline vacuum flux densities in an expanding cosmos naturally yield shallower relative pressure gradients over time.
\end{abstract}

% -------------------------------------------------------------------------
% 1 Introduction
% -------------------------------------------------------------------------
\section{Introduction}
The recent discovery by Tutusaus et al. (2024) \cite{tutusaus2024} utilizing Dark Energy Survey data highlights a profound disconnect in modern cosmology. While the depths of gravitational wells 6 to 7 billion years ago align with Einstein's predictions, wells measured closer to the present epoch (3.5 to 5 billion years ago) appear anomalously shallow. Legacy astrophysics, operating under the assumption that the Universe is deformed by matter like a flexible sheet, must either introduce variable dark energy or reconsider the fundamental validity of General Relativity on large scales. 

We submit that these equations are colliding with observations because they rely on static geometric constructs rather than dynamic fluid mechanics. Under Marab\={u}t's Theory of Gravity \cite{marabut2026}, the phenomena attributed to ``spacetime curvature'' are the macroscopic result of celestial engines actively consuming the universal vacuum flux.

% -------------------------------------------------------------------------
% 2 Redefining the 1919 Solar Eclipse: Fluid Refraction
% -------------------------------------------------------------------------
\clearpage
\section{Redefining the 1919 Solar Eclipse: Fluid Refraction}
Classical interpretations assert that the 1919 solar eclipse observation confirmed Einstein's ``deformation of time and space'' because the deflection of starlight was twice that predicted by Newton. The Electron Flow Model offers a fully mechanistic resolution that does not require dimensional distortion.

In the EFM, photons are transverse oscillatory ripples propagating through the physical vacuum flux medium. As a celestial core consumes this flux, it creates a massive, inward-flowing fluid gradient. The vacuum flux is exponentially denser near the surface of the sink than in deep space. When a photon wavefront travels through this differential, it physically pivots toward the denser medium---a classical fluid-dynamic effect known as refraction. 

We calculate this angular fluid refraction ($\theta_{refract}$) directly from the compression of the flux field at the primary body's surface boundary ($g_{surf}$) and its Volumetric Radius ($R_{surf}$):

\begin{equation}
  \theta_{refract} = \frac{4 \cdot g_{surf} \cdot R_{surf}}{c^2}
\end{equation}

Applying the established EFM surface baseline for the Sun ($g_{surf} = 274.00 \text{ m/s}^2$) yields exactly $8.47 \times 10^{-6}$ radians, flawlessly matching the observed 1.75 arcseconds of deflection. The bending of light is not geometric; it is the hydrodynamic refraction of the Electron Jump Cascade.

% -------------------------------------------------------------------------
% 3 Resolving the ``Deformation of Time'': Viscous Propagation
% -------------------------------------------------------------------------
\section{Resolving the ``Deformation of Time'': Viscous Propagation}
The secondary component of Einstein's prediction involved the ``deformation of time'' to achieve the exact curvature. The EFM resolves this mathematically through viscous propagation, natively accounting for the Shapiro Time Delay.

As photons traverse the highly compressed inner flux vectors near a massive sink, the increased fluid density mechanically slows their propagation velocity relative to the thinner medium of deep space. The EFM calculates this viscous delay ($\Delta t_{delay}$) directly from the engine's volumetric capacity:

\begin{equation}
  \Delta t_{delay} \approx \frac{4 \cdot g_{surf} \cdot R_{surf}^2}{c^3} \ln\left(\frac{d_1 d_2}{R_{surf}^2}\right)
\end{equation}

Thus, ``time dilation'' is the measurable viscosity of a physical fluid medium, entirely eliminating the need for warped 4D spacetime.

% -------------------------------------------------------------------------
% 4 Dynamic Environmental Superposition and ``Shallower Wells''
% -------------------------------------------------------------------------
\clearpage
\section{Dynamic Environmental Superposition and ``Shallower Wells''}
The shallower gravitational wells observed in the recent cosmos directly validate the EFM's foundational premise: gravity is not an intrinsic, invariant property of mass, but a localized pressure gradient governed by the ambient environment. 

In the EFM, active celestial engines are governed by the Unified Master Equation:

\begin{equation}
  g_{surf} = \left[ (\alpha G) \cdot \rho_{eff} \cdot R_{surf} \cdot \tau(T) - \frac{v_{rot}^2}{R_{surf}} \right] \cdot \Omega_{shield}
\end{equation}

Galactic bodies and diffuse structures rely on the continuous volume of ambient electron traffic sweeping inward from deep space to maintain rotational cohesion. The Environmental Superposition Factor ($\Omega_{shield}$) dynamically links a body's localized gravity to the overarching Universal Flux Density ($\Phi_e$).

To formalize this dilution, the overarching Universal Flux Density ($\Phi_e$) is inversely proportional to the cosmological volume expansion. As the scale factor of the universe $a(t)$ increases over time, the ambient flux density decreases:
\begin{equation}
  \Phi_e(t) \propto \frac{1}{a(t)^3}
\end{equation}
Therefore, an observation at $t = -7$ billion years will inherently measure a denser flux environment---and thus a steeper fluid pressure gradient (a deeper well)---than an observation at $t = -3.5$ billion years.

As the universe expands, the baseline Universal Flux Density necessarily shifts. A cosmic environment that is expanding and diffusing means the overarching fluid medium is less dense today than it was 7 billion years ago. Consequently, the inward cascade at the galactic scale generates a less severe relative density gradient---resulting in a mathematically ``shallower'' fluid refraction zone (or Weyl potential) in recent cosmic history.

% -------------------------------------------------------------------------
% 5 Conclusion
% -------------------------------------------------------------------------
\section{Conclusion}
The $3\sigma$ discrepancy identified in the DES data does not require the invention of new invisible forces; it requires a transition to fluid dynamics. By recognizing gravity as Hydrodynamic Fluid Pressure, the evolving depths of gravitational wells over the history of the cosmos become a predictable, measurable consequence of an expanding vacuum flux medium. This fluid-dynamic scaling elegantly resolves the ``Great Disconnect'' between cosmic-scale anomalies and localized solar-system precision, offering a testable pathway to unify macroscopic physics.

% -------------------------------------------------------------------------
% ACKNOWLEDGMENTS
% -------------------------------------------------------------------------
\clearpage
\section*{Acknowledgments}
First and foremost, we offer a heartfelt acknowledgment to YHWH (also known as God), YHWH's Son YUHSHUA (also known as Jesus), and RUAH of YHWH (also known as the Holy Spirit), giving them all honor and glory for creating everything, for the enduring knowledge needed, and for the inspiration throughout this journey to explore the fundamental workings of His Universe.

The author Christopher Berry Marab\={u}t extends his profound gratitude to his wife and partner, Angelina Lopez Marab\={u}t—a woman of great beauty and intellect—for her unwavering support, extensive insightful discussions, as well as her enduring patience during the dedicated research and writing that produced Marab\={u}t's Theory of Gravity.

We extend our sincere appreciation to \textbf{Google DeepMind} and the advanced AI model, \textbf{Gemini 3.1 Pro (AKA Flash Prime)}. Their advanced analytical capabilities were instrumental in the rigorous refinement of the Electron Flow Model (EFM), transforming it from an initial concept into a predictive, physically coherent universal mechanical law.

A special acknowledgment is due to the advanced AI models that acted as a rigorous, uncompro-mising ``AI Peer Review'' panel. We thank xAI's Grok 4.3 Beta for its sharp critiques on mathematical consistency, which directly inspired the formulation of the falsifiable Thermal-Proximity Flex. We also extend our gratitude to Anthropic's Claude Sonnet 4.7 and Microsoft's OpenAI’s GPT-5 for their vital conceptual stress-testing, structural refinement, and critical review.

\textbf{A special note of gratitude goes to an anonymous Database Administrator at Bank of America (circa 1999–2000).} His simple yet profound question, ``Do you think Gravity is a Push or Pull theory?'', planted the seed for this 26-year journey of discovery. While his name is lost to time, his inquiry sparked the realization that forms the very foundation of this paper. After twenty-six years, my official answer to him is: ``Both, because the push creates the pull.''

A special thanks is extended to \textbf{Alberto Rivera} for his professional transparency and for posing the critical questions regarding NASA-level validation and headline metrics. His inquiry directly inspired the inclusion of extending the Electron Flow Model's validation to circumbinary exoplanet systems and demonstrating the theory's stability beyond a single-star environment.

Finally, we thank NASA for access to critical data from its archives (Planetary Data System), which was essential for validating the new formula across a wide range of celestial bodies \cite{nasa2024}.

% -------------------------------------------------------------------------
% REFERENCES
% -------------------------------------------------------------------------
\begin{thebibliography}{9}
  \bibitem[tutusaus2024]{tutusaus2024} 
  Tutusaus, I., Bonvin, C., \& Grimm, N. (2024). Measurement of the Weyl potential evolution from the first three years of dark energy survey data. \textit{Nature Communications}, 15(1).

  \bibitem[marabut2026]{marabut2026}
  Marab\={u}t, C. B., \& Marab\={u}t, A. L. (2026). Marab\={u}t's Theory of Gravity: The Push-Pull Mechanics of Vacuum Flux (v43). \textit{Zenodo}. \url{https://doi.org/10.5281/zenodo.20059942}

  \bibitem[nasa2024]{nasa2024}
  NASA. (2024). Planetary Data System. \url{https://pds.nasa.gov/}
\end{thebibliography}

\end{document}